\chapter{คลังโค้ด Python สำหรับ OpenXR STEM}
\label{app:C}

\section*{ค.1 Template พื้นฐาน OpenXR Session Loop}

\begin{codebox}[OpenXR Session Loop Template ภาษา Python]
import openxr as xr
import numpy as np

def create_openxr_session():
    """สร้าง OpenXR Instance, System, และ Session"""
    # สร้าง Instance
    instance_info = xr.InstanceCreateInfo(
        application_info=xr.ApplicationInfo(
            application_name="STEM Lab",
            application_version=xr.MAKE_VERSION(1, 0, 0),
            engine_name="RBRU Physics Engine",
            engine_version=xr.MAKE_VERSION(1, 0, 0),
            api_version=xr.CURRENT_API_VERSION,
        ),
    )
    instance = xr.create_instance(instance_info)
    
    # เลือก System (HMD)
    system_id = xr.get_system(
        instance,
        xr.SystemGetInfo(form_factor=xr.FormFactor.HEAD_MOUNTED_DISPLAY)
    )
    return instance, system_id

def run_frame_loop(session, action_set):
    """วงรอบ Frame หลักสำหรับการเรนเดอร์ XR"""
    while True:
        # Poll system events
        event = xr.poll_event(session.instance)
        if event is None:
            break
        
        # รอสถานะ Running
        frame_state = xr.wait_frame(session)
        xr.begin_frame(session)
        
        if frame_state.should_render:
            # --- Render STEM virtual lab here ---
            render_virtual_lab(session, frame_state.predicted_display_time)
        
        xr.end_frame(session, xr.FrameEndInfo(
            display_time=frame_state.predicted_display_time,
            environment_blend_mode=xr.EnvironmentBlendMode.OPAQUE,
        ))
\end{codebox}

\section*{ค.2 การคำนวณ Pinch Distance สำหรับ Hand Interaction}

\begin{codebox}[คำนวณ Pinch Gesture จาก 21-Joint Hand Data]
def compute_pinch_distance(thumb_tip: np.ndarray, 
                            index_tip: np.ndarray) -> float:
    """คำนวณระยะระหว่างปลายนิ้วโป้งกับปลายนิ้วชี้"""
    return float(np.linalg.norm(thumb_tip - index_tip))

def is_pinching(joints: list, threshold: float = 0.025) -> bool:
    """ตรวจสอบว่ากำลังจีบนิ้ว (Pinch) หรือไม่"""
    # XR_HAND_JOINT_THUMB_TIP_EXT = 5
    # XR_HAND_JOINT_INDEX_TIP_EXT = 10
    thumb_pos = np.array(joints[5].pose.position)
    index_pos = np.array(joints[10].pose.position)
    dist = compute_pinch_distance(thumb_pos, index_pos)
    return dist < threshold

def detect_gestures(joints: list) -> dict:
    """ตรวจจับท่าทางมือ: Pinch, Grasp, Point"""
    return {
        "pinch": is_pinching(joints, threshold=0.025),
        "point": is_pointing(joints),
        "grasp": is_grasping(joints),
    }
\end{codebox}

\section*{ค.3 Physics Integrator สำหรับ Virtual STEM Lab}

\begin{codebox}[RK4 Integrator สำหรับการจำลองฟิสิกส์เสมือน]
import numpy as np
from dataclasses import dataclass

@dataclass
class RigidBody:
    mass: float
    position: np.ndarray   # เวกเตอร์ตำแหน่ง [x, y, z]
    velocity: np.ndarray   # เวกเตอร์ความเร็ว [vx, vy, vz]
    force: np.ndarray      # แรงรวมที่กระทำ [Fx, Fy, Fz]

def rk4_step(body: RigidBody, dt: float) -> RigidBody:
    """อินทิเกรต 1 ก้าวด้วย Runge-Kutta อันดับ 4"""
    def deriv(pos, vel):
        acc = body.force / body.mass
        return vel, acc
    
    k1_v, k1_a = deriv(body.position, body.velocity)
    k2_v, k2_a = deriv(body.position + 0.5*dt*k1_v, body.velocity + 0.5*dt*k1_a)
    k3_v, k3_a = deriv(body.position + 0.5*dt*k2_v, body.velocity + 0.5*dt*k2_a)
    k4_v, k4_a = deriv(body.position + dt*k3_v, body.velocity + dt*k3_a)
    
    new_pos = body.position + (dt/6.0)*(k1_v + 2*k2_v + 2*k3_v + k4_v)
    new_vel = body.velocity + (dt/6.0)*(k1_a + 2*k2_a + 2*k3_a + k4_a)
    
    return RigidBody(body.mass, new_pos, new_vel, body.force)
\end{codebox}
